% Section VI - A reporting checklist. Table II items are transcribed VERBATIM from blueprint Sec 11
% (advisor content). Do not reword or add items; Sec 14: keep to eight, answerable in about ten minutes.
% Budget: ~220 words of prose plus the table (word_budget_MLESP.md).

% REVIEW 2026-11-07 (8-page limit): the three prose paragraphs condensed into one. No item changed.
Table~\ref{tab:checklist} lists eight items a connectivity-based study should state, one per leakage surface, each answerable in a sentence; together they let a reader decide from the paper alone whether test information could have reached training, a question \PctUnclearAny\% of the audited studies left open for at least one step. For items 5 to 7 what matters is which data a step was fit on, and the epochs per subject (item 2) show how much a subject-identity shortcut could contribute under epoch-level partitioning. The checklist adds no procedure: a study that already partitions by subject and fits every statistic within the training fold satisfies it with a few sentences. Each item should be answered even when the step was not performed, so a reader can tell an omitted step from an undescribed one---a distinction this audit could not draw (Sections~\ref{sec:methods},~\ref{sec:results}).

% REVIEW 2026-11-07 (R2 W11, Q7): relation to existing reporting guidance.
% The TRIPOD+AI and CLAIM entries were REMOVED from references.bib: at the 8-page limit their two reference
% lines were the overflow onto page 9, and blueprint Sec 12 forbids citing a paper no author has opened, which
% there was no time to do. They are named, not cited; only brookshire2024, already verified, carries a citation.
General reporting guidance for prediction models and medical-imaging AI---TRIPOD+AI, CLAIM---already asks for the validation scheme and the handling of predictors, and recommendations for EEG decoding ask for subject-level partitioning \cite{brookshire2024}, but none names the statistics specific to a connectivity matrix. Table~\ref{tab:checklist} replaces none of them; it adds that layer.

\begin{table}[t]
\caption{Reporting checklist for connectivity-based deep learning studies (Gen.\ = generalization)}
\label{tab:checklist}
\centering
\small
\begin{tabular}{@{}r p{0.70\columnwidth} l@{}}
\toprule
\# & A connectivity study should state & Addresses \\
\midrule
1 & Whether data were partitioned by subject or by epoch & L1 \\
2 & The number of subjects and the number of epochs per subject & L1 \\
3 & Whether connectivity was estimated separately within each training fold & L2 \\
4 & The window length and overlap used for connectivity estimation & L2 \\
5 & How any threshold or edge-count parameter was selected, and from which data & L3 \\
6 & Whether feature normalization statistics were computed on training data only & L4 \\
7 & Whether channel or band selection used test-set performance & L5 \\
8 & Whether any external or cross-dataset evaluation was performed & Gen. \\
\bottomrule
\end{tabular}
\end{table}
